% compass-aqsgd-latex-additions.tex
%
% The paper delta for the AQ-SGD codec ablation, kept in the repo because the
% Overleaf source lives outside version control and two sessions were editing
% it concurrently. If the paper file loses these edits, re-apply from here.
%
% Target: main_rewritten.tex in the COMPASS Overleaf project. Red markup uses
% the project's existing \new{...} macro and revised environment.
%
% THREE INLINE INSERTIONS (anchor text -> what to append after it):
%
% 1. Related Work, activation-compression paragraph. After:
%      "...require compatible model structure \citep{ramasinghe2025subspace,
%       ramasinghe2025mixtures,wang2021pufferfish}."
%    append:
%      \new{Our RLVR codec selection compares against AQ-SGD directly rather
%      than only citing it, at matched boundary payload;
%      Appendix~\ref{app:rlvr-aqsgd} reports the arm and the reason its premise
%      is a poor fit for on-policy RLVR.}
%
% 2. Table 6 (tab:rlvr-codec-choice) gains a second block. The caption is
%    rewritten and five rows are added. See the table body committed alongside
%    this file in the launcher's run plan, or reconstruct from Table 6 in the
%    compiled PDF.
%
% 3. Conclusion and Limitations. After:
%      "...only one pair runs for $600$ steps."
%    append:
%      \new{The codec comparison in Appendix~\ref{app:rlvr-aqsgd} is
%      single-seed per arm, so a difference smaller than run-to-run variation
%      cannot be distinguished from one, and it is measured at one payload
%      point on one pair.}
%
% PLUS the appendix section below, appended before \end{document}.
%
% Verified: adds ZERO new LaTeX errors (204 pre-existing tectonic
% \resizebox+tabular artifacts both before and after), zero unresolved
% references, renders as Appendix K over 5 subsections.

\begin{revised}
\section{The Quantization Family at Matched Payload}
\label{app:rlvr-aqsgd}

Table~\ref{tab:rlvr-codec-choice} selects a codec by comparing a low-rank
family against a sparsifying one. Activation quantization is the third
established family for this link, and AQ-SGD \citep{wang2022aqsgd} is its
reference method for pipeline parallelism over slow networks. This appendix
reports it at matched boundary payload on Pair~1, together with the two control
arms needed to attribute the result.

\subsection{What AQ-SGD Sends}

AQ-SGD does not quantize the activation. It quantizes the \emph{change} of the
activation for the same training example between visits. Each side of a
boundary keeps a local buffer $m(\xi)$ for example $\xi$. The sending stage
transmits $Q(a(\xi)-m(\xi))$ and both sides advance
\begin{equation}
m(\xi)\leftarrow m(\xi)+Q\big(a(\xi)-m(\xi)\big),
\label{eq:aqsgd-update}
\end{equation}
so the receiving stage consumes a running reconstruction. The backward gradient
is quantized directly, with no buffer, which is the paper's own treatment: a
gradient carries no cross-visit history to difference against. The guarantee
rests on a self-enforcing loop. Training stabilizes, so the activation for a
fixed example moves less between visits, so the same bit budget spent on the
delta gives a smaller error, so training stabilizes further.

Our implementation sends a keyed subset. For each token a pseudorandom function
selects exactly $k$ coordinates, using the same key construction as the mask
codec so that both adjacent stages derive the subset locally and no index bits
cross the link. Unlike the mask, the coordinates outside the subset are not
zeroed and carry no $H/k$ gain, because the receiver already holds $m$ there and
keeps it. The reconstruction is
\begin{equation}
\hat h \;=\; m \;+\; \mathcal S_J\!\left(Q\big(h_J-m_J\big)\right),
\label{eq:aqsgd-subset}
\end{equation}
where $\mathcal S_J$ scatters back onto the selected coordinates $J$. Over
successive visits this refines $m$ toward $h$ one subset at a time. Rounding is
stochastic, which is the faithful setting: Theorem 3.1 of
\citet{wang2022aqsgd} assumes an unbiased quantization function, and the
worked bound in that paper rounds stochastically. What the result does not
assume is an unbiased \emph{gradient}, which is a different claim. A
round-to-nearest arm is reported alongside to show that this choice is not what
decided the outcome.

\subsection{Payload, and Why the Arms Are Comparable}

At $H=1536$ the mask sends $77$ half-precision coordinates, so $1232$ bits per
token per boundary. Quantization at $b$ bits on $k$ selected coordinates sends
$kb$ payload bits plus one half-precision scale per block of $32$ selected
coordinates, so $kb + 16k/32$ bits. At $b=2$ and $k=493$ that is $1232.5$ bits,
matching the mask to within $0.04\%$ while covering $32.1\%$ of each token
rather than $5.0\%$. The memoryless arm uses identical $(b,k,\text{block})$
values, so it is byte-identical to the AQ-SGD arms on the wire and differs in
exactly one respect, whether a buffered history is differenced against. Any gap
between those arms is therefore attributable to delta coding alone.

One setting cannot be made byte-matched, and we report it separately rather
than adjust it away. AQ-SGD sends the first message for each example
uncompressed. Where most tokens are first visits, that dominates the payload
and the method does not reach a $5\%$ budget at all. The rows in
Table~\ref{tab:rlvr-codec-choice} therefore use a cold fallback that respects
the budget, in which a first visit sends the unbiased sparse estimate and is
consequently identical to the memoryless arm. A separate short run with the
faithful uncompressed first message measures the resulting payload directly.

\subsection{Why the Premise Does Not Transfer}

AQ-SGD needs a training example to recur with its activation intact. On-policy
RLVR only half provides this, and the split is what the two AQ-SGD rows
measure.

Prompt-prefix tokens do recur. A MATH problem is revisited each epoch with
byte-identical token ids at identical positions, so its boundary activation
differs only by the weight drift accumulated since the last visit. At $600$
steps and $128$ prompts per step over the $7{,}498$ training problems, a prompt
is seen $10.2$ times. This is the regime AQ-SGD was designed for.

Response tokens do not recur. The token occupying a given position is resampled
from the current policy at every visit, so a buffered value belongs to a
different token. For roughly independent activations
$\mathbb E\|a-m\|^2=\|a\|^2+\|m\|^2$, so the delta carries about $\sqrt2$ times
the norm of the value. At a fixed bit budget that is a coarser grid and a
larger error than quantizing the value would have given. We measure this ratio
at $1.41$ on Pair~1 boundary states.

The prefix is also a minority of what crosses a boundary, and we measure that
share rather than estimate it. Tokenized with the run's own chat template, the
$7{,}498$ MATH prompts have a mean length of $111.4$ tokens (median $82$, $95$th
percentile $280$), against a mean response length of $574$. The prefix is
therefore $16.3\%$ of the tokens that cross a boundary, which is a ceiling on
how much of the traffic AQ-SGD's mechanism can reach at all, before any
question of whether it helps there. The arm labelled ``recurring prefix''
confines the buffer to those positions and is AQ-SGD's best case here. The arm
labelled ``all positions'' is the unmodified port. We report the measured
buffer hit rate with both, so the outcome is explained rather than asserted.

\subsection{Storage}

The buffer is indexed by example, position and boundary, and it must span the
reuse distance to score a hit. In RLVR that distance is a full epoch of
prompts. The cost is $n\cdot\bar p\cdot L\cdot H\cdot 2$ bytes for $n$ examples
of mean prompt length $\bar p$ over $L$ compressed boundaries, which at Pair~1
is $0.150$~GiB per buffered token position. Covering prefixes alone therefore
needs $16.7$~GiB of host memory, and covering full sequences to the $2048$-token
budget needs $308$~GiB. Pair~3's $16{,}384$-token budget puts the same quantity
in the terabytes. We run the prefix arm with a $20$~GiB cap, so it never evicts
inside an epoch and the hit rate it reports is not an artifact of the cap; the
all-positions arm runs under the same cap and does evict, which is itself part
of the measurement. The keyed mask stores nothing at
all, because its selection is derived from a key rather than remembered. Buffers
are also per process, which is faithful: \citet{wang2022aqsgd} note that data
parallelism divides the store by the parallel degree but then pays for
shuffling, so a prompt that lands on a different replica than last epoch is a
miss. This is a property of the method in this setting and not an artifact of
our implementation, and it is the reason we report a hit rate alongside every
accuracy.

\subsection{Protocol}

All arms use the published Pair~1 configuration with the boundary codec
replaced and nothing else changed: Qwen2.5-Math-1.5B on MATH, $1024$ prompt and
$2048$ response tokens, $128$ prompts per step with $G=8$ and a $128$-prompt
mini-batch so that one optimizer tick follows each generation, anchor cadence
and delay $20/20$, and $600$ optimizer steps. Each arm runs on a single GPU,
matching the published rows: at data-parallel degree above one the anchor's
dense replay gradient is averaged across replicas, which changes the quantity
under study. Final checkpoints are merged to dense models and scored with the
same ten-benchmark harness that produced Table~\ref{tab:rlvr-full}, so the new
rows are read against the Pair~1 dense and mask controls already reported there
rather than against the selection run. Checkpoints are retained at steps $200$,
$400$ and $600$, and training-path MATH validation runs every $50$ steps,
because a single end-of-run reading cannot distinguish a codec that trains
steadily from one that peaks and then drifts.
\end{revised}
